% !TeX root = ../papers/nc-p3.tex
\section{極大イデアルの二乗が単項となる平面局所点}
\label{sec:finite-stretched-planar-marks}
\begin{theorem}\label{thm:finite-stretched-planar-marks}
四辺が有限な反対階数8の標準有限Frobenius窓に、接次元2で極大イデアルの二乗が
単項の局所点が一つあれば、全次数の標準正則性は有限に決まる。
従ってこの階数の有限枝で既出の正則性判定が被覆していない候補は、
閉点が高々三つで、全ての局所環が接次元2、長さ6以上、極大イデアルの二乗が非単項の場合に限る。
\end{theorem}
\begin{proof}[局所環と三重点切断]
三文字認証がないと仮定し、当該点を印とする。局所環 $P$ はGorensteinであり、
パラメータを取れば
\[
 P=k[x,y]/(xy,\ y^2-x^e,\ x^{e+1}),\qquad e\ge2 .
\]
実際 $\mathfrak m^2/\mathfrak m^3$ は一次元なので標数零で一般の $x$ の平方がその生成元となる。
$e>2$ なら $\mathfrak m^3/\mathfrak m^4$ も一次元で、一般の $x$ を $x^3\ne0\bmod\mathfrak m^4$ とも
選べる（三次の偏極による）。このとき $\mathfrak m^3=(x^3)$ で、順に $\mathfrak m^j=(x^j)$ ($j\ge2$)。
実際 $\mathfrak m x^j=x^{j-2}\mathfrak m x^2$ である。$e=2$ でも以下同様である。
第二生成元から $x$ の多項式を引いて $xy=0$ とする。すると $y^2\in\mathfrak m^2$
は $x$ に消されるので $x^e$ の倍で、零ならsocleの次元に反する。倍率を揃えて表示を得る。
どの $f\in\mathfrak m$ に対しても、非零の商 $P/\operatorname{ann}f$ は曲線状である。
実際 $\mathfrak m f$ は一列のイデアル $(x^j)$ からなる $\mathfrak m^2$ の部分加群だからである。
$P/(y)=k[x]/x^e$ とし、各元を $x$ の次数 $e$ 以下の多項式と $y$ の定数倍に分ける。
$\mathfrak m$ の三元以上の積の $x$ 表示は、各 $y$ 項を捨てて通常に掛け法 $x^{e+1}$ で取ったものに等しい。

共通切断の $\Sigma$ が正次元または長さ2以下なら既出の平面点判定を使う。
長さ3のGorenstein型での印の評価誤差を考える。正規化した二行列の差を $Z$ とする。
その全成分は $\mathfrak m$ に入る。定理\ref{thm:finite-marked-artin-secants}で使った三方向の
スキームで消える二次式 $Q$ について、$Q(l) Z=0$ を示せば以後の証明は変わらない。
$y$ を捨てた $Z$ の成分の最小 $x$ 次数を $s\ge1$ とし、$s\ge e-1$ なら三因子の積より直ちに従う。
その他では $k[x]/x^e$ において $Z=x^s V'$ とし $l,r$ を $x$ で割る。
商での $l,r$ の座標は極大イデアルを生成するので、割った列の剰余は非零。
$A,V'$ は法 $x^{e-s-1}$ で独立な相対割線の基底となる（正規化によりその剰余は独立）。
従って $Q(l/x)=0\bmod x^{e-s-1}$。二次の $x^2$ と $x^s$ を戻すと少なくとも
$e+1$ 次となる。三因子の積の表示により元の環でも零である。

非Gorenstein三重点では定理\ref{thm:finite-marked-planar-triple}の $t,l_0$ と評価 $f$ を使う。
印では $f,t,l_0\in\mathfrak m$ である。曲線状商 $P/(y)$ での元の割線の証明は
$t f$ がそのsocleに、$l_0 f$ が零になることを与える。元の環でこれらは $\mathfrak m^2$
に入るので、各々 $(x^{e-1}),(x^e)$ に含まれる（$e=2$ でも長さの上界は同じ）。
印以外で零だからそれぞれ長さ2以下で、対応する隣の一次座標がこの支持環に全射する。
隣接座標積の自己直交性より両方とも零である。これで前定理の長さ6の完全交叉 $Y$ と、
誤差全体を生成する一元 $f$ の議論が適用できる。
印の $P/\operatorname{ann} f$ 自身が曲線状なので、$Z=fV'$ を書いた後、
この環で $l,r$ を一様化変数で割れば、局所長さを一つ減らした上で $A,V'$ の相対割線を得る。
$A$ と $V'$ は五つの共有関係の評価核にあり（後者は差を $f$ で割る）、
ともに割線の平面 $t=0$ に入る。従って $Y$ の長さ5以上で余因子が消えるという矛盾まで同じである。
\end{proof}

\begin{proof}[長さ四の切断での二次の誤差]
定理\ref{thm:finite-artin-kernel-symmetry}と同じ外積 $k$、その成分イデアル $(\eta)$、
支持環 $S=B/I$ を用いる。今は印も $Z=\operatorname{Spec}S$ に入ってよい。
$W'$ が全射の場合と共線四点因子の場合の排除には印の仮定がない。
残る余次元1での分母と、印以外での乗法行列の議論もそのまま成り立つ。

$n=3$ のとき $f=H\chi(a)$ と書く。印で
\[
                          \mathfrak m^2 f\subset(\eta)
\]
を示す。$(\eta)$ を法とし、さらに $H$ の零化イデアルで割る。
この環上で三行の乗法行列 $M=m_L$ は全二次小行列式が零で $m_{\bar a}M=\alpha I_3$。
従って $\alpha M=0$ である。実際 $\alpha M_{ij}$ は $k\ne j$ として
$\sum_h(m_{\bar a})_{kh}(M_{hk}M_{ij}-M_{hj}M_{ik})$ に等しい。
$M$ の成分は $l$ の三座標を張る（$V\to S$ は同型）。
$I$ を先頭の基底とする $m_l$ の下三行は $(0,M)$ である。
それらの二次小行列式が零で、$m_a m_l=\alpha I_4$ の第一行を用いれば
$\chi(a)$ は $m_l$ の二次小行列式を全て消す。実際第一行を含む小行列式に
$\chi(a)$ を掛け、第一行を積の第一行で置き換えると、残りの項は下三行の小行列式、
または $\alpha$ と下三行の成分との積である。

これらの小行列式は印で $\mathfrak m^2$ を生成する。接方向の二次元鉛筆は
$V$ の二平面（$l$ が接ベクトルを分離する）で、どの非零方向でも乗法階数は2以上である。
その二次小行列式が全二項二次式を張ることを確認する。さもなくばそれらを消す
二次汎関数を取り、平方ならある方向の階数が1以下になる。
相異なる二因子なら基底変更後の二行列 $U,V'$ の混合二次小行列式が全て零である。
$U$ を階数 $r\ge2$ の単位ブロックにして微分すると $V'$ はそのブロックの外で零となる
（非対角成分には別の非零対角元を使う）。$r\ge3$ ならブロック内でも非対角は零、
対角はどの二つの和も零なので $V'=0$。$r=2$ なら二行二列の鉛筆となり、
行列式に根がある。いずれも矛盾する。次数環へ写しNakayamaを使えば生成が従う。
先の二つの商から戻れば表示の包含を得る。

$B$ が局所 $(1,2,1)$ で $n=4$ の場合にも $f=H\chi(a)$ に同じ包含が成り立つ。
同じ二つの商で $l=s+v$、$v$ の二つの接座標を $v_1,v_2$ とすると、逆数の分子より
$s^2=s v_1=s v_2=0$。$al=\alpha 1$ の一次部分は $\chi(a)v_i+s a_i=0$ なので
$\chi(a)v_i v_j=0$ でもある。$s,v_1,v_2$ は $l$ の三座標の基底である。

どちらでも $f$ が印で単元で $(\eta)$ が零なら $\mathfrak m^2=0$ で矛盾する。
印で $(\eta)\ne0$ なら、$a(p)$ は $S$ の支持に属する指標の跡双対元である。
$n=3$ の $\chi$ は欠けた長さ1のイデアルの指標で、この双対元に非零となるのは
その指標の因子が単純で同じ指標の場合だけである。その場合当該因子は $S$ に残っていない。
局所 $(1,2,1)$ の場合もsocle元の剰余は零。従って常に $f$ は印で非単元であり、
その局所支持環は曲線状である。よって $(f)/\mathfrak m^2(f)$ の長さは高々2であり、
全体の支持長は $n+2$ 以下となる。$n=3$ は下界6に反し、$n=4$ で等号なら
$(\eta)\subset(f)$ が全点で成り立つため支持が $Z$ を含み、偶座標の平面条件に反する。
\end{proof}

\begin{proof}[曲線状因子の積からなる長さ四の切断]
$n=4$ で $B$ が曲線状因子の積の場合は $\operatorname{adj}m_l$ の各成分が $(\eta)$ に入り、
印以外の $\operatorname{adj}m_a$ の消滅は既に証明した。印では $a(p)$ の乗法階数は1である。
実際 $A(p)$ は接次元2により階数1で $\Sigma$ に属する。
従って $m_a$ の三次小行列式は $\mathfrak m^2$ に入る。$e\le4$ ではこれは長さ3以下である。
$e\ge5$ とする。$\eta$ は外積から $\mathfrak m^2$ に入り、
$\eta\ne0$ なら $(\eta)=(x^s)$、$n_p=e+1-s$ とする。$\eta=0$ なら $s=e+1,n_p=0$ とする。
$Z$ 上で $m_a$ が階数1なので、その二次小行列式の $x$ 部分は法 $x^{n_p}$ で零である。
$s\le3$ ならこれだけで三次小行列式には $x^{e-2}$ 以上の次数がある。
$s>3$ では、$l$ の三因子の積の公式から、$y$ を捨てた $l$ を $x$ で割った $l'$ の
乗法の三次小行列式は法 $x^{s-3}$ で零である。この割算は $s=e+1$ でも、
三因子の積を法 $x^{e+1}$ で読めることから正当である。$l'$ の乗法には単元な二次小行列式がある。
法 $x^{e-1}$ で $al'=\alpha'1$ と書け、さらに法 $x^{s-3}$ では像が階数2なので
前定理の生成元数の議論で $\alpha'=0$。従って法 $x^{e-1}$ では、二つの基底列を $l'B$ から選んで
$m_a$ の三次小行列式を展開すると全て $x^{s-3} I_2(m_a)$ に入る。
$s-3+n_p=e-2$ だから、元の環でも三次小行列式は長さ3の $(x^{e-2})$ に入る
（既に $\mathfrak m^2$ に入るため $y$ の一次項はない）。
独立な四つの三次式をここに入れることはできず、長さ4の場合も核側の対称性を得た。

片側対称性からの正則性判定は系\ref{cor:all-defective-one-sym-finite}による。
中央支持が全て4なら印で $w(p_0,p_2)\ne0$ であり、接次元は3以下の行列の余階数として高々2である。
接次元2の長さ4または5の局所Gorenstein環は、平方零の二次元socleを持てないため
$\dim\mathfrak m^2/\mathfrak m^3=1$ である（長さ5でこれが2なら $\mathfrak m^3=0$）。
長さ20の和を取れば残余の閉点数の上界を得る。
\end{proof}
